\documentclass[10pt]{article}
\usepackage[margin=0.75in]{geometry}
\usepackage{amsmath,amssymb,mathtools}
\usepackage{mathptmx}
\DeclareMathAlphabet{\mathcal}{OMS}{cmsy}{m}{n}
\usepackage{booktabs}
\usepackage{enumitem}
\usepackage[hidelinks]{hyperref}
\usepackage{titlesec}
\setlist{nosep,leftmargin=*}
\titlespacing*{\section}{0pt}{8pt}{3pt}
\titlespacing*{\subsection}{0pt}{5pt}{2pt}
\setlength{\parskip}{3pt}
\setlength{\parindent}{0pt}
\newcommand{\B}{\mathsf{B}}
\newcommand{\Rd}{\mathsf{R}}
\newcommand{\As}{\mathrm{A}}
\newcommand{\Pt}{\mathrm{P}}
\newcommand{\1}{\mathbf{1}}

\title{\vspace{-2.5em}\textbf{Workshop 2 plan: Bayesian beliefs and a multi-period best response}\\[2pt]
\large Proposal for Sean, Nov 1, 2026 (nothing built yet)}
\date{\vspace{-3em}}

\begin{document}
\maketitle

\section*{Summary}
\begin{enumerate}
\item \textbf{Bayesian beliefs.} Blue does not know how often Red asserts when Blue ships are present. It holds a
Beta prior on Red's assertion rate in each region and updates it with Bayes' rule each period. It sees Red's
posture only in zones where it has ships, so it learns only where it goes. That is exactly the setting from class
where a Bayesian decision maker can get stuck on a wrong belief forever (a self-confirming equilibrium). We compare
four ways to choose under the belief (myopic, a committed presence patrol, Thompson sampling, an optimistic bound) and
measure each by regret against a Blue that knows Red's true rates.
\item \textbf{Multi-period best response.} Today's MIP-BR plans one period ahead and only considers staying, going
home, or moving one edge. The new version plans $H$ periods ahead on a time-expanded copy of the map, carries out
the first period, and solves again next period. It lets Blue position ships for zones two or more moves away and
plan around endurance.
\item Both use only methods from class. Bayes' rule with a Beta prior, self-confirming equilibrium, committed
experiments, Thompson sampling, optimism and regret come from the learning material. Types and posterior beliefs
come from incomplete-information games. Discounting comes from dynamic games, and the time-expanded flows are the
network-flow MIP from the optimization material.
\end{enumerate}

\section{Bayesian beliefs about Red's behavior}

\subsection{What Blue does not know, and what it sees}
Red's \emph{type} is its assertion rate in each region $r\in\{N,C,S\}$,
\[
\theta_r=\Pr(\text{Red asserts at a zone in region } r \mid \text{Blue ships are present there}),\qquad
\theta=(\theta_N,\theta_C,\theta_S)\in[0,1]^3 .
\]
Red knows $\theta$. Blue does not. Blue's prior is $\mu_0(\theta_r)=\mathrm{Beta}(\alpha_r,\beta_r)$, independent
across regions.

\textbf{One model change is needed.} Today a detection at $v$ (own ships, adjacent ships, or satellite) also
reveals Red's posture at $v$. I propose that counts can still be seen from a distance, but posture (whether Red is
shouldering, shadowing or warning) is seen only where Blue has ships, which matches the fog-of-war rule
``only forces present where they also maintain forces''. This is the feedback constraint from class. Blue sees the
outcome only of the zones it chooses to go to, and nothing about the zones it stays out of.

\subsection{Bayes update}
In period $t$, let $k_{r,t}$ be the number of Blue-occupied zones in region $r$ where Blue saw Red assert, and
$m_{r,t}$ the number where it saw Red patrol. Then
\[
\mu_{t+1}(\theta_r)=\mathrm{Beta}\Big(\alpha_r+\textstyle\sum_{s\le t}k_{r,s},\ \beta_r+\sum_{s\le t}m_{r,s}\Big),
\qquad
\hat q_{r,t}=\frac{\alpha_r+\sum_{s\le t}k_{r,s}}{\alpha_r+\beta_r+\sum_{s\le t}(k_{r,s}+m_{r,s})}.
\]
A region Blue stays out of adds nothing, so its posterior does not move.

\subsection{How the belief enters the best response}
Today the option values $a_v(n,c,p)$ and risks $\rho_v(n,c,p)$ assume Red's last-seen posture at $v$ (patrol if
never seen). With a belief they become an expectation over Red's two postures, with $q$ the rate the choice rule
uses in region $r(v)$:
\[
\bar a_v(n,c,p)=(1-q)\,a_v(n,c,p\mid \Pt)+q\,a_v(n,c,p\mid \As),\qquad
\bar\rho_v(n,c,p)=(1-q)\,\rho_v(n,c,p\mid \Pt)+q\,\rho_v(n,c,p\mid \As).
\]
These are computed before the solve, so the MIP keeps exactly its current form and size.

\subsection{Four choice rules}
\begin{center}\small
\begin{tabular}{lll}
\toprule
Rule & Rate $q$ used for region $r$ in period $t$ & What it does\\\midrule
Myopic & $\hat q_{r,t}$ (posterior mean) & acts on the current belief, never explores on purpose\\
Committed presence patrol & $\hat q_{r,t}$, plus at least one ship in every region for the first $N_e$ periods & a fixed experiment, then myopic\\
Thompson sampling & a fresh draw $\tilde\theta_{r,t}\sim\mu_t(\theta_r)$ & explores in proportion to what it does not yet know\\
Optimistic bound & $\max\!\big\{0,\ h_r/n_r-\sqrt{2\ln t/n_r}\big\}$ & judges each region by the lowest rate not yet ruled out\\
Oracle (benchmark) & the true $\theta_r$ & the full-information best response\\
\bottomrule
\end{tabular}\end{center}
$h_r$ and $n_r$ are the assertions seen and zone-periods observed in region $r$ so far. For Blue a low assertion
rate is good news, so optimism subtracts the bonus. The committed patrol is a realistic option for the US side
(a scheduled presence operation regardless of what Red does). Thompson draws use their own random stream, so games
stay on common random numbers.

\subsection{The trap in this game}
If Blue starts pessimistic about a region (for example South with prior mean $0.8$), the expected value of
going there is low and the expected risk is high, so the myopic MIP keeps Blue out. Blue then sees nothing in the
South, its belief never moves, and Red holds the South for the rest of the year. Nothing Blue observes contradicts
its belief, so this is a self-confirming equilibrium. Blue is optimal given its prediction, and its predictions are
right everywhere it actually goes. The questions are how often the myopic rule gets trapped, how much each other
rule pays to avoid it, and whether design D2 (extra surveillance) breaks the trap by reporting Red's posture in zones
next to Blue ships. That is the ``reports from other ships'' counter from class. It also gives D2 a concrete reason
to matter, since it showed no benefit so far.

\subsection{Regret}
For a fixed true $\theta$, regret over the year is
\[
R_T(\theta)=\mathbb E\Big[\sum_{t=1}^{T}\big(J^\theta_t(x^\star_t)-J^\theta_t(x_t)\big)\Big],\qquad T=180,
\]
where $J^\theta_t$ is the MIP objective at the period-$t$ state with option values computed from the true $\theta$,
$x_t$ is the learner's plan, and $x^\star_t$ is the oracle's plan for the same state. Each Blue decision needs one
extra solve (the oracle's), the same check we ran for greedy against the MIP. We also report the game-level gap
$u_\B(\text{oracle})-u_\B(\text{learner})$ on the same seeds.

\subsection{Which Red Blue learns against}
\begin{itemize}
\item \textbf{Correctly specified.} Red moves its ships with MIP-BR, but where Blue is present it asserts with
probability $\theta_r$ each period. Blue's model contains the truth, so a learner that keeps exploring should reach
it. Two type profiles, for example $\theta=(0.2,0.5,0.8)$ and $(0.8,0.5,0.2)$.
\item \textbf{Misspecified.} Red is the normal MIP-BR, whose posture depends on the state and on Blue. No single
$\theta$ describes it, so we report where Blue's beliefs settle and what that costs, which is the ``wrong model''
failure from class.
\end{itemize}
Red learning about Blue at the same time would make each side's environment change as the other learns. Class has
not covered learning by two sides yet, so I would leave that out unless it is taught before Dec 1.

\subsection{Optional second piece: where unseen ships are}
Old sightings still count fully today. A Bayes-consistent fix: if each Red ship leaves its node with probability $m$
per period, a sighting of age $a$ is still right with probability $(1-m)^a$. The believed count blends the last
sighting with the unseen ships spread over the zones, with weight $(1-m)^a$. From the Workshop 1 runs $m\approx0.4$.
This is small and makes D2's fresher sightings worth something. I would add it only if time allows after the above.

\section{Multi-period best response}

\subsection{Why}
The one-period MIP only looks at staying, going home, or moving one edge. A zone two or more moves away has no
value to it until a ship happens to be next to it, it cannot plan a two-step move into the South, and it does not
see that a ship sent out now will have to turn back for fuel in three periods. These are the effects the grader is
most likely to question (``one-period lookahead'').

\subsection{Formulation}
Make $H$ copies of the 17-node map, one per period $\tau=1,\dots,H$. Each ship type moves along the edges of this
time-expanded graph (stay, or one edge per period), which is a min-cost flow over time with integer flows. Let
$n_{v,\tau}$ and $c_{v,\tau}$ be the navy and coast guard levels at zone $v$ in period $\tau$, and $p_{v,\tau}$ the
planned posture. Then
\[
\max\ \sum_{\tau=1}^{H}\delta^{\tau-1}\Big(\sum_{v}\bar a_v(n_{v,\tau},c_{v,\tau},p_{v,\tau})-\gamma\,M_\tau\Big)
\quad\text{s.t.}\quad
M_\tau\le F_i,\qquad \sum_v\bar\rho_v(n_{v,\tau},c_{v,\tau},p_{v,\tau})\le\varepsilon_i\Big(1-\frac{\ell}{L}\Big)+\xi_\tau
\ \ \text{for each }\tau,
\]
plus flow conservation, endurance (a ship must stay within reach of a port or the oiler at every step), and the same
level and option constraints as today in every period. $M_\tau$ counts the moves in period $\tau$. Blue carries out
period 1 only and solves again next period with new observations.

\textbf{What it assumes.} Over the $H$ periods Red's ships and posture are held at Blue's current belief (with the
decay above, if used). It does not model Red reacting to Blue's plan inside the lookahead, so it is still a best
response to a forecast, not a solution of the dynamic game. Holding the belief fixed also means the lookahead never
values a move for what it would teach Blue, so exploration still has to come from the choice rule. Whether looking
ahead can do that job is the next part of the learning material.

\subsection{Cost}
The one-period MIP has at most 339 variables (measured over a 20-period MIP-BR game). With $H$ periods it grows
roughly $H$ times, so $H=2$ fits the size-limited license and $H=3$ may not. Sean's academic license, or a Gurobi
module on Engaging if one exists (\texttt{module avail gurobi}), removes the limit. Solve time grows faster than
size, so I would measure $H=2$ and $H=3$ on a few games before picking.

\section{Experiments and timing}
\begin{center}\small
\begin{tabular}{lll}
\toprule
Block & Setups & Games (200 seeds each)\\\midrule
L1 learning, correct model & 4 rules $\times$ 2 priors $\times$ 2 Red types, plus oracle & 3{,}600\\
L2 learning, misspecified & 4 rules against MIP-BR Red & 800\\
L3 trap and D2 & myopic, pessimistic prior, D0 and D2 & 400\\
M1 lookahead & $H\in\{2,3\}$ Blue against rule-based, MIP-BR, equilibrium Red & 1{,}200\\
\bottomrule
\end{tabular}\end{center}
About 6{,}000 new games, run on Engaging the same way as the suite (one overnight). Learning games cost about one
extra solve per Blue decision. Lookahead games cost more per solve.

\textbf{Order.} Beliefs, the four rules and regret by Nov 10, then a short local test and the Engaging run before the
Nov 17 learning checkpoint. Multi-period MIP by Nov 24, then its runs. Writeup and slides after Dec 1, with QA.

\section*{Decisions for Sean}
\begin{enumerate}
\item Posture seen only where Blue has ships (counts still seen from a distance)?
\item Learn per region (3 rates, faster learning) or per zone (8 rates)?
\item Blue learns, Red has a fixed type. Or should Red also learn?
\item Lookahead $H=2$ or $H=3$, and is it worth asking ORCD about a Gurobi license on Engaging?
\item Include the optional unseen-ships belief?
\end{enumerate}
\end{document}
